\chapter{AI Models: Moisture Prediction and Process Novelty}
\label{chap:models}

\section{Two different questions, two different models}

A soft sensor estimates a difficult or slow-to-obtain quality variable from easier measurements. This is a common process-industry role, especially when laboratory data arrive slowly, but it also inherits problems such as multi-rate sampling, collinearity, missing data and drift \cite{kadlec2009}. In this project, the supervised branch answers: \emph{what final moisture is expected from the available context?}

The second branch asks a different question: \emph{does the current process vector resemble the reference regime?} A One-Class Support Vector Machine (One-Class SVM) learns a boundary around reference observations rather than named failure classes \cite{scholkopf2001}. Its output cannot explain a cause by itself. I therefore added a separate diagnostic layer that ranks deviations and suggests what should be checked.

\section{Soft sensor: governed experiment}

Notebook 03 is the source of truth for model selection and the runtime artifact. It receives 1,103 causally aligned rows with 16 ordered inputs: nine direct process values, five engineered interactions and the strictly previous laboratory density and product temperature. The final moisture at the current laboratory timestamp is the target. The chronological split is 772 TRAIN, 166 VALIDATION and 165 untouched TEST rows.

I evaluate the candidates with mean absolute error (MAE), root mean squared error (RMSE) and the coefficient of determination ($R^2$). Together, these metrics show average error, sensitivity to larger errors and explained variation.

Within TRAIN, five expanding \code{TimeSeriesSplit} folds tune each candidate. The later VALIDATION block selects the model family; TEST is opened once after that choice. This preserves chronology and avoids the optimistic mixing that random cross-validation can create in dependent series \cite{bergmeir2012}.

\begin{table}[htbp]
\centering\small
\caption{Audited contract of the selected moisture model}
\label{tab:moisture-provenance}
\begin{tabularx}{\textwidth}{@{}p{0.29\textwidth}Y@{}}
\toprule
Item & Audited value \\
\midrule
Source & \code{MAP\_Dryer\_Lab\_Aligned\_16.csv}; 1,103 laboratory-aligned rows. \\
Chronology & TRAIN 772; VALIDATION 166; TEST 165; no shuffle. \\
Selection & Expanding five-fold tuning inside TRAIN; family selected by later VALIDATION RMSE. \\
Pipeline & Standardization followed by Ridge Regression, $\alpha=10$. \\
Contract & Exactly 16 original ordered inputs; final moisture excluded from features. \\
Artifact & \code{models/5s/quality\_moisture\_pipeline.joblib}. \\
\bottomrule
\end{tabularx}
\source{Executed Notebook 03 and saved evaluation record \cite{moisturenotebook}.}
\end{table}

Five model families were compared under the same structure. Ridge obtained the lowest VALIDATION RMSE, 0.0013566 percentage points, narrowly ahead of Elastic Net at 0.0013625 and ordinary Linear Regression at 0.0013671. The Ridge/Elastic Net difference is only about 0.43\%, so I treat Ridge as the selected version for this dataset and split, not as a universal winner.

\begin{figure}[htbp]
\centering
\includegraphics[width=.90\textwidth]{notebook03_candidate_validation.png}
\caption{Later-segment VALIDATION RMSE used after training-only tuning.}
\source{Executed Notebook 03 \cite{moisturenotebook}.}
\label{fig:moisture-candidates}
\end{figure}

\section{Holdout performance and baseline value}

After selection, Ridge is refitted on TRAIN plus VALIDATION and evaluated on the 165-sample TEST tail. It obtains MAE 0.0010686, RMSE 0.0014026, $R^2=0.8245295$, bias 0.0000662 and maximum absolute error 0.0056206 percentage points H$_2$O.

A model score is more informative when compared with a simple operational alternative. I therefore added a persistence baseline that predicts the current moisture using the strictly previous laboratory moisture. On the same TEST rows, persistence gives MAE 0.0014589 and RMSE 0.0018479. Ridge reduces those errors by 26.8\% and 24.1\%, respectively.

\begin{table}[htbp]
\centering\small
\caption{Untouched TEST comparison with the last-laboratory persistence baseline}
\label{tab:baseline-comparison}
\begin{tabular}{@{}lrrrrr@{}}
\toprule
Method & $n$ & MAE & RMSE & $R^2$ & Bias \\
\midrule
Previous lab moisture & 165 & 0.001459 & 0.001848 & 0.695439 & 0.000012 \\
Selected Ridge & 165 & \textbf{0.001069} & \textbf{0.001403} & \textbf{0.824530} & 0.000066 \\
\bottomrule
\end{tabular}
\source{Independent report audit using the frozen artifact and identical TEST timestamps.}
\end{table}

\begin{figure}[p]
\centering
\includegraphics[width=.95\textwidth]{notebook03_test_holdout.png}
\caption{Measured and predicted moisture on the untouched chronological TEST interval.}
\source{Executed Notebook 03 \cite{moisturenotebook}.}
\label{fig:notebook-moisture-evidence}
\end{figure}

\begin{figure}[p]
\centering
\includegraphics[width=.96\textwidth]{soft_sensor_baseline_residuals.png}
\caption{Persistence comparison and residual audit for the selected Ridge soft sensor.}
\source{Independent report evidence generated from the frozen artifact.}
\label{fig:softsensor-residuals}
\end{figure}

The residual mean is 0.0000662 and the median is 0.0000550; 47.9\% of residuals are negative and 52.1\% positive. The 5th and 95th percentiles are $-0.002170$ and 0.002152. A Durbin--Watson statistic of 1.827 and lag-1 residual autocorrelation of 0.084 show limited remaining short-lag structure in this synthetic TEST sequence. Absolute residual and fitted value have correlation 0.054, so no strong level-dependent error pattern is visible. These diagnostics are useful, but 165 sparse targets from one synthetic chronology remain too little evidence for plant acceptance.

Coefficient magnitude is also interpreted cautiously. Air per feed and wet-product feed have the largest absolute standardized coefficients, followed by cooler- and dryer-air temperature. The inputs are correlated, so weight can be redistributed among them; a large Ridge coefficient is not a causal effect.

\section{One-Class SVM for earlier process awareness}

Dryer variables evolve together, and an unusual combination can appear before the next laboratory result. I use the One-Class SVM as a process-novelty detector: it compares the current multivariate state with the learned reference region. A novel state prompts investigation; it is not a fault probability, a root-cause classification or a control command.

Notebook 04 exports a radial basis function (RBF) One-Class SVM with \code{gamma='scale'} and $\nu=0.02$ \cite{anomalynotebook}. It fits 6,000 sampled rows from the canonical TRAIN period and checks a separate 20,000-row validation sample. The fraction of validation states identified as novel is 5.925\%, reported here as the \emph{validation novelty flag rate}.

The raw decision function is a signed model score, not a probability. The service converts it to a zero-to-one display risk for easier reading, with 0.5 at the detector boundary and 0.8 as the visual high-risk level. These values support prioritization on the dashboard; they do not represent failure probability or confidence.

\begin{figure}[htbp]
\centering
\begin{tikzpicture}[node distance=4mm and 5mm,scale=.82,transform shape]
\node[data,text width=23mm] (state) {Process\\state};
\node[flow,text width=25mm,right=of state] (score) {One-Class SVM\\novelty score};
\node[flow,text width=24mm,right=of score] (severity) {Anomaly state\\and severity};
\node[flow,text width=25mm,right=of severity] (rank) {Ranked\\contributors};
\node[data,text width=27mm,right=of rank] (subsystem) {Likely subsystem\\possible causes};
\node[data,text width=28mm,right=of subsystem] (verify) {Recommended check\\and operator interpretation};
\draw[arr] (state)--(score); \draw[arr] (score)--(severity); \draw[arr] (severity)--(rank);
\draw[arr] (rank)--(subsystem); \draw[arr] (subsystem)--(verify);
\end{tikzpicture}
\caption{Implemented path from process novelty to operator verification.}
\source{Author's synthesis of the runtime and diagnostic workflow.}
\label{fig:novelty-workflow}
\end{figure}

\section{Held-out process-novelty response}

The chronological TEST interval contains 237,600 rows, including 1,620 rows from three labelled synthetic disturbance episodes. The frozen detector flags 6,305 rows in total, or 2.654\% of the interval. It flags 60.25\% of rows inside the labelled episodes and 2.258\% of rows outside them. These are scenario-response and normal-period flag rates, not accuracy or plant fault probabilities.

\begin{table}[htbp]
\centering\small
\caption{Frozen novelty-detector response to the three labelled TEST episodes}
\label{tab:scenario-response}
\begin{tabularx}{\textwidth}{@{}lrrrY@{}}
\toprule
Episode & Rows & Flagged & Row response & First flag / interpretation \\
\midrule
Feed surge & 408 & 0 & 0.0\% & No flag; peak display risk 0.360. \\
Steam dip & 540 & 443 & 82.0\% & First raw flag after 290 s; peak risk 0.978. \\
Airflow restriction & 672 & 533 & 79.3\% & First raw flag after 360 s; peak risk 0.950. \\
\bottomrule
\end{tabularx}
\source{Independent report audit on the held-out canonical TEST interval.}
\end{table}

The detector reacts strongly to the steam-dip and airflow-restriction departures, although sensitivity varies with disturbance type and magnitude. It does not flag the feed-surge episode. This is useful design feedback: representative operating modes and better loading-related features matter more than presenting the model as a complete fault classifier.

\section{From novelty to investigation guidance}

The useful contribution is the full diagnostic chain, not only the binary flag. The runtime compares a novel row with the TRAIN-period reference profile, ranks direct-variable deviations, gathers likely-subsystem evidence and presents possible causes with recommended verification. Detection, evidence localization and causal diagnosis remain separate tasks \cite{venkatasubramanian2003}.

In a representative dashboard demonstration event, vacuum, air flow and fan speed are the strongest contributors. Their joint movement points to the drying-air and exhaust subsystem, so the interface suggests checking fan demand, damper position, leakage and the vacuum measurement. This is a probable diagnosis for investigation, not a confirmed root cause.

The detector was developed and evaluated on deterministic synthetic process data and should therefore be interpreted as a proof of concept for process-novelty monitoring. Its operating threshold and sensitivity would need to be recalibrated using representative plant history, operating modes and confirmed events before industrial use.

\begin{engineeringbox}[How the anomaly branch is used]
The detector acts as an early-warning layer. It highlights a process state, localizes the strongest evidence and leaves verification and action to the operator using PCS 7, laboratory information and normal operating procedures.
\end{engineeringbox}

\begin{figure}[htbp]
\centering
\includegraphics[width=.94\textwidth]{anomaly_holdout_validation.png}
\caption{One-Class Support Vector Machine (OCSVM) decision function and display risk on the chronological TEST replay.}
\source{Frozen One-Class SVM evaluated without refitting.}
\label{fig:anomaly-holdout}
\end{figure}

\FloatBarrier
